\section{``Solving'' means unexploitable, and unexploitable is a number}\label{sec:solving}

\mkey{what solving means}Machine learning usually optimizes against a fixed target: a loss goes down, accuracy goes up. Poker has no fixed target. The quality of any strategy depends entirely on what the opponent does, and the opponent adapts. So the field's definition of ``solved'' is borrowed from game theory: a strategy pair is a \keyterm{Nash equilibrium}\mdef{Nash equilibrium}{a pair of strategies where neither player can gain by changing only their own} when neither player can improve by deviating unilaterally. \hlight{An equilibrium strategy is one you could publish --- an opponent who reads your exact strategy still cannot beat it in expectation.}

Three classical facts make this the right target for \emph{heads-up} poker, and they are worth internalizing before any algorithm:

\begin{enumerate}
\item \textbf{Two-player zero-sum is special.} Heads-up poker is \keyterm{zero-sum}\mdef{zero-sum}{whatever one player wins, the other loses; payoffs sum to a constant} --- every chip you win, the opponent loses. Von Neumann's minimax theorem says such games have a unique \emph{value} (the expected win rate of perfect play), and all equilibria achieve it. Better: equilibria are \emph{interchangeable} --- any equilibrium strategy of yours is safe against any strategy of theirs, guaranteeing at least the game value. This is exactly what fails with three or more players, which is why every guarantee in this survey is heads-up-only (Section~\ref{sec:checklist} returns to this).
\item \textbf{Equilibrium play must be random.} If you only bet when strong, your bets stop getting paid; if you never bluff, your checks get attacked. The equilibrium bluffs, calls, and value-bets with calibrated \emph{frequencies} --- it is a \keyterm{mixed strategy}\mdef{mixed strategy}{a probability distribution over actions, rather than one fixed action}. \pitfall{This single fact --- optimal play is deliberately stochastic --- is what breaks most standard RL machinery later in this survey}: greedy ``take the best action'' improvement has no fixed point to find.
\item \textbf{Distance from equilibrium is measurable.} Fix your strategy $\sigma$ and let a perfect adversary --- the \keyterm{best response}\mdef{best response}{the strategy maximizing EV against one known, fixed opponent strategy} --- extract the maximum from it. What it wins above the game value is your \keyterm{exploitability}\mdef{exploitability}{how much a perfect counter-strategy beats your strategy for; zero exactly at equilibrium}. Exploitability is the loss function of this entire field: every algorithm in this survey is judged by how fast it drives that number toward zero, and Section~\ref{sec:evaluation} is about how to measure it at each scale.
\end{enumerate}

\subsection*{The object being solved: a game tree with hidden cards}

\mkey{extensive form}Poker is played on a tree: chance nodes deal cards, decision nodes belong to a player, leaves pay out the pot. The formal name is an \emph{extensive-form game}. Hidden information enters through one definition. Suppose you hold $K\,K\,9\,5\,2$ after a bet; there are many tree nodes consistent with what you see --- one per opponent hand --- and you cannot tell them apart. That set of indistinguishable nodes is one \keyterm{information set}\mdef{information set}{all game states a player cannot tell apart: own cards + public cards + the action history} (``infoset''), and a strategy must assign \emph{one} action distribution per infoset, not per node. Everything downstream --- what CFR tables index, what networks take as input, what the dashboard displays --- is keyed by infosets.

\begin{intuition}{hidden information means you play against a probability cloud}
In chess you respond to \emph{the} position. In poker you respond to a cloud: every opponent hand, weighted by how likely it is to have played this way --- their \emph{range}. An equilibrium strategy is one whose action frequencies hold up against the whole cloud at once. This is also why poker strategies are ranges, not hands: the dashboard question is never ``what does he have?'' but ``what fraction of everything he could have does this?''
\end{intuition}

One technical condition matters enough to name: \keyterm{perfect recall}\mdef{perfect recall}{no player ever forgets their own past observations or actions} --- players remember everything they saw and did. Kuhn's theorem (1953) guarantees that under perfect recall, per-infoset randomization (``behavioral strategies'') loses no generality, and CFR's convergence proof quietly depends on it. It becomes a live issue exactly once in this survey: ``imperfect-recall abstraction'' (Section~\ref{sec:abstraction}) deliberately violates it to save memory, trading the guarantee away.

\mnote{Drawmaha's split pot changes payoffs at the leaves only --- the game stays zero-sum, and every guarantee in this section applies unchanged.}Heads-up Drawmaha fits this frame with no modification. The split pot makes leaf payoffs more intricate (halves, scoops), but payoffs still sum to the pot: the game is constant-sum, equivalently zero-sum, and the minimax/interchangeability/exploitability story carries over exactly.
